\documentclass[tikz,border=2pt]{standalone}
\usepackage{amsmath}
\usepackage{pgfplots}
\pgfplotsset{compat=1.18}

\begin{document}
% The zeros of f are the inputs where the graph meets the x-axis: f(x)=x^3-3x+1 has three.
\begin{tikzpicture}
	\begin{axis}[
		axis lines=middle,
		xmin=-2.6, xmax=2.6,
		ymin=-2.4, ymax=3.6,
		xtick={-1.879,0.347,1.532}, xticklabels={,,},
		ytick=\empty,
		xlabel={$x$}, ylabel={$y$},
		xlabel style={below right}, ylabel style={above left},
		samples=150, smooth, clip=false,
		width=8cm, height=5.6cm
		]
		\addplot[thick, blue, domain=-2.25:2.1] {x^3-3*x+1};
		\addplot[only marks, mark=*, red] coordinates {(-1.879,0) (0.347,0) (1.532,0)};
		\node[red, anchor=north west, font=\small] at (axis cs:-1.82,-0.08) {$x_1$};
		\node[red, anchor=south west, font=\small] at (axis cs:0.4,0.05) {$x_2$};
		\node[red, anchor=north west, font=\small] at (axis cs:1.6,-0.08) {$x_3$};
		\node[blue, anchor=west] at (axis cs:2.1,3.2) {$f$};
	\end{axis}
\end{tikzpicture}
\end{document}
